%% ========================================================================

%% Bhaskara II: Bijaganita Part II — Bilingual Sanskrit-English Edition

%% LEFT COLUMN: Sanskrit (Devanagari)

%% RIGHT COLUMN: English Translation with IAST

%% ========================================================================

\documentclass[11pt,a4paper]{article}



%% -------- Core Packages --------

\usepackage{fontspec}

\usepackage{polyglossia}

\usepackage{amsmath,amssymb,amsthm}

\usepackage{geometry}

\usepackage{graphicx}

\usepackage{xcolor}

\usepackage{titlesec}

\usepackage{fancyhdr}

\usepackage{enumitem}

\usepackage{array,booktabs,longtable,multirow}

\usepackage{hyperref}

\usepackage{lettrine}

% \usepackage{microtype}  % disabled: incompatible with ucharclasses font switching

\usepackage{tocloft}

\usepackage{indentfirst}

\usepackage{ucharclasses}

\usepackage{tikz}

\usepackage{float}

\usepackage{caption}

\usepackage{subcaption}



%% -------- Page Geometry (wider margins for parallel columns) --------

\geometry{

  paperwidth=210mm,paperheight=297mm,

  left=18mm,right=18mm,top=25mm,bottom=25mm,

  headheight=14pt,footskip=30pt

}



%% -------- Language Setup --------

\setmainlanguage{english}

\setotherlanguage{sanskrit}



%% -------- Font Configuration --------

\setmainfont{Times New Roman}

\setsansfont{Arial}

\setmonofont{Courier New}



%% Devanagari Font

\newfontfamily\devanagarifont[Script=Devanagari,Scale=1.05]{Nirmala UI}

\newfontfamily\devanagarifontsf[Script=Devanagari,Scale=1.05]{Nirmala UI}



%% ucharclasses: automatic font switching for mixed scripts

\setDefaultTransitions{\normalfont}{}

\setTransitionsForDevanagari{\devanagarifont}{}



%% -------- Colors --------

\definecolor{titlecolor}{RGB}{45,55,72}

\definecolor{sectioncolor}{RGB}{55,65,81}

\definecolor{notecolor}{RGB}{120,53,15}

\definecolor{rulecolor}{RGB}{180,160,140}

\definecolor{versecolor}{RGB}{60,50,40}

\definecolor{sanskritbg}{RGB}{255,253,245}

\definecolor{englishbg}{RGB}{250,252,255}

\definecolor{columnrule}{RGB}{160,140,120}



%% -------- Section Formatting --------

\titleformat{\section}

  {\normalfont\Large\bfseries\color{sectioncolor}}

  {\thesection}{1em}{}

\titleformat{\subsection}

  {\normalfont\large\bfseries\color{sectioncolor}}

  {\thesubsection}{1em}{}



%% -------- Header/Footer --------

\pagestyle{fancy}

\fancyhf{}

\fancyhead[L]{\small\textit{B\={\i}jaga\d{n}ita — Bilingual Edition}}

\fancyhead[R]{\small\thepage}

\renewcommand{\headrulewidth}{0.4pt}



%% -------- Custom Commands --------

\newcommand{\longline}{\noindent\rule{\textwidth}{0.4pt}}

\newcommand{\shortline}{\noindent\rule{0.5\textwidth}{0.4pt}\par}

\newcommand{\cnote}[1]{\textcolor{notecolor}{\textit{#1}}}

\newcommand{\dev}[1]{{\devanagarifont #1}}

\newcommand{\vs}{\vspace{0.5em}}



%% Sanskrit verse environment for bilingual display

\newenvironment{sanskritverse}{%

  \begin{quote}\devanagarifont\textcolor{versecolor}%

}{%

  \end{quote}

}



%% Column environment: Sanskrit | English

\newcommand{\sktcol}[1]{%

  \noindent\begin{minipage}[t]{0.47\textwidth}%

  \setlength{\parindent}{0pt}%

  #1%

  \end{minipage}%

}

\newcommand{\engcol}[1]{%

  \begin{minipage}[t]{0.47\textwidth}%

  \setlength{\parindent}{0pt}%

  \raggedright

  #1%

  \end{minipage}%

}

\newcommand{\sktenvpair}[2]{%

  \par\vspace{0.6em}%

  \noindent\sktcol{#1}\hfill\vrule\hfill\engcol{#1}%

  \par\vspace{0.6em}%

}



%% Bilingual pair command: \bilingual{Sanskrit}{English}

%% ucharclasses handles automatic Devanagari font switching

\newcommand{\bilingual}[2]{%

  \par\vspace{0.5em}%

  \noindent\begin{minipage}[t]{0.47\textwidth}%

  \setlength{\parindent}{0pt}%

  \small #1%

  \end{minipage}%

  \hfill\textcolor{columnrule}{\vrule width 0.5pt}\hfill%

  \begin{minipage}[t]{0.47\textwidth}%

  \setlength{\parindent}{0pt}%

  \small\raggedright #2%

  \end{minipage}%

  \par\vspace{0.5em}%

}



%% Verse numbering

\newcounter{verseno}

\newcommand{\verseref}[1]{\hfill\textcolor{notecolor}{\textsc{(#1)}}}



%% Shloka macro for bilingual display

\newcommand{\shlokaSkt}[1]{%

  \begin{center}\textcolor{versecolor}{\devanagarifont #1}\end{center}%

}

\newcommand{\shlokaEng}[1]{%

  \begin{center}\textit{#1}\end{center}%

}



%% -------- Hyperref Setup --------

\hypersetup{

  colorlinks=true,

  linkcolor=sectioncolor,

  citecolor=sectioncolor,

  urlcolor=sectioncolor,

  pdfencoding=auto

}



%% -------- TOC Depth --------

\setcounter{tocdepth}{2}



%% ========================================================================

\begin{document}



%% -------- Title Page --------

\begin{titlepage}

\centering

\vspace*{1.5cm}



{\Huge\bfseries\color{titlecolor} B\={\i}jaga\d{n}ita\par}

\vspace{0.3cm}

{\Large\itshape The Algebra of Bh\={a}skara II\par}



\vspace{1cm}

{\LARGE\bfseries\color{sectioncolor} Part II\par}

\vspace{0.2cm}

{\large Pages 58--114\par}



\vspace{1.5cm}



\begin{center}

\rule{0.6\textwidth}{1pt}

\end{center}



\vspace{0.3cm}

{\large\devanagarifont

बीजगणितम् — द्विभाषीसंस्करणम्\\[0.3em]

Bilingual Sanskrit-English Edition\par}



\vspace{0.3cm}

\begin{center}

\rule{0.6\textwidth}{1pt}

\end{center}



\vspace{1cm}



{\large\devanagarifont

अथ एकवर्णसमीकरणम्\\

अव्यक्तवर्गादिसमीकरणम्\\

अनेकवर्णसमीकरणम्\\[0.5em]

अनेकवर्णमध्यमाहरणभेदाः\par}



\vspace{0.5cm}

{\normalsize

1.~Equations with One Unknown (Ekavar\d{n}a-sam\={\i}kara\d{n}a)\\

2.~Equations with Unknown Squares (Avyakta-varg\={a}di-sam\={\i}kara\d{n}a)\\

3.~Equations with Many Unknowns (Anekavar\d{n}a-sam\={\i}kara\d{n}a)\\

4.~Variations of the Elimination Method (Anekavar\d{n}a-madhyam\={a}hara\d{n}a-bhed\={a}h)\par}



\vfill



{\large Bh\={a}skara II (1114--1185 CE)\par}

\vspace{0.3cm}

{\small Sanskrit text in Devanagari $\parallel$ English translation with IAST\par}



\end{titlepage}



%% -------- Table of Contents --------

\tableofcontents

\newpage



%% ========================================================================

\section*{Introduction / प्रस्तावना}

\addcontentsline{toc}{section}{Introduction}



\bilingual{\textbf{संस्कृतम्}\\[0.5em]

अयं ग्रन्थः श्रीभास्कराचार्यविरचितस्य \textbf{बीजगणितस्य} द्वितीयं भागं समाविशति। बीजगणितं गणितशास्त्रस्य मूलभूतं खण्डं यत् समीकरणानां साधनं, अज्ञातराशीनां निर्णयं च व्याख्यायते। पृष्ठक्रमांक ५८--११४ पर्यन्तं चत्वारि प्रधानाध्यायानि समाविष्टानि—एकवर्णसमीकरणम्, अव्यक्तवर्गादिसमीकरणम्, अनेकवर्णसमीकरणम्, तथा अनेकवर्णमध्यमाहरणभेदाः।

}{%

\textbf{English Translation}\\[0.5em]

This volume contains \textbf{Part II} of the \textit{B\={\i}jaga\d{n}ita} (``Seed Computation,'' i.e., Algebra) composed by \textbf{Bh\={a}skara II} (1114--1185 CE). The \textit{B\={\i}jaga\d{n}ita} is the foundational treatise of Indian algebra, treating equations, unknown quantities, and their solutions. Pages 58--114 encompass four principal chapters: equations with one unknown, equations with unknown squares, equations with many unknowns, and variations of the elimination method.}



\vs



\bilingual{भास्कराचार्यस्य गणितस्य विशेषता अस्ति यत् प्रत्येकं विषयं श्लोकेन (सूत्रेण) उपदिश्य, तदनन्तरं उदाहरणैः (worked examples) सह व्याख्याति। अज्ञातराशिः \dev{यावत्तावत्} (``as much as'') इत्युच्यते, संक्षेपेण \dev{या} अथवा \dev{याव} इति लिख्यते। समीकरणानि \textit{समशोधन} (elimination/equation) क्रियया साध्यन्ते।

}{%

A distinctive feature of Bh\={a}skara's exposition is that each topic is introduced by a verse (\textit{s\={u}tra}) containing a rule, followed by worked examples (\textit{ud\={a}hara\d{n}a}) and detailed commentary (\textit{v\={a}rttika}) showing the method of solution step by step. The unknown quantity is denoted by \dev{यावत्तावत्} (\textit{y\={a}vatt\={a}vat}, ``as much as''), abbreviated \dev{या} or \dev{याव}. Equations are solved by the process of \textit{sama-\'{s}odhana} (making equal and clearing).}



\vs



\bilingual{\textbf{मुख्याः पद्धतयः} — अस्य भागस्य प्रधानगणितपद्धतयः सन्ति: (१) \textbf{चक्रवालपद्धतिः} (cyclic method) — \textit{वर्गप्रकृतेः} ($x^2 - Ny^2 = 1$) साधनार्थं भास्करस्य श्रेष्ठमल्गोरिद्मिकसिद्धान्तः; (२) \textbf{भावितपद्धतिः} — घातसमीकरणानां साधनार्थं; (३) \textbf{मध्यमाहरणम्} — बहुवर्णसमीकरणानां साधनार्थं निरसनपद्धतिः।

}{%

\textbf{Key Methods} treated in this part include: (1) The \textbf{Cakrav\={a}la} (cyclic) \textbf{method} --- Bh\={a}skara's greatest algorithmic achievement for solving \textit{vargaprak\d{r}ti} ($x^2 - Ny^2 = 1$, Pell's equation); (2) The \textbf{Bh\={a}vita} method for product equations; (3) \textbf{Madhyam\={a}hara\d{n}a} --- the elimination method for equations with multiple unknowns.}



\newpage



%% ========================================================================

\part*{Chapter I: Equations with One Unknown}

\addcontentsline{toc}{part}{Chapter I: Equations with One Unknown}



\section*{\dev{एकवर्णसमीकरणम्} $\parallel$ Ekavar\d{n}a-sam\={\i}kara\d{n}a}

\addcontentsline{toc}{section}{Ekavar\d{n}a-sam\={\i}kara\d{n}a (pp.~58--76)}



\bilingual{\lettrine[lines=2, lraise=0.1, nindent=0em, slope=-.5em]{\dev{अ}}{}थैकवर्णसमीकरणम्। अत्र \textbf{यावत्तावत्} एको वर्णः प्रकल्प्यते, ततो रूपैः सह समीकृत्य यावत्तावन्मानं साध्यते। राशिश्च द्विधा—\textit{स्वीकृतपक्षः} (accepting side) तथा \textit{दातृपक्षः} (giving side) इति। तयोः समशोधनेन मानं निर्णीयते।

}{%

\lettrine[lines=2, lraise=0.1, nindent=0em, slope=-.5em]{N}{}ow, equations with one unknown. Here, a single unknown quantity, called \textbf{y\={a}vatt\={a}vat} (\dev{यावत्तावत्}, ``as much as''), is postulated. Combining it with known quantities (\textit{r\={u}pa}, \dev{रूप}), one forms an equation and determines the value of the unknown. Each side of the equation is called either the \textit{accepting side} (what is received) or the \textit{giving side} (what is given). By the process of equal-clearing (\textit{sama-\'{s}odhana}), the value is determined.}



\vs



\subsection*{Simple Linear Equation $\parallel$ सरलरेखीयसमीकरणम्}



\bilingual{\textbf{उदाहरणम्} — अत्राश्वमूल्यमज्ञातं तस्य मानं \dev{यावत्तावत्} एकं प्रकल्प्तं \dev{या १}। तदा श्रेणीशकं यदेकस्य यावत्तावन्मूल्यं तदा पणां किमिति कल्पितच्छागपणां प्रमाणमकं लब्धं पणामश्वानां मूल्यम्। \dev{या ६}। अत्र रूपातनये प्रक्षिप्ते जातमाश्चस्य धनं \dev{या ६ रूप २०}। एवं दशानां मूल्यं \dev{या १०}। अत्र रूपाते चाणांते प्रक्षिप्ते जातं द्वितीयस्य धनं \dev{या १० रूप १००}।

}{%

\textbf{Example} --- Here, the price of a horse is unknown; let its value be one \textbf{y\={a}vatt\={a}vat}, written \textbf{y\={a} 1}. If the price of one horse equals the price of so many goats, then the price in coins is determined. For one horse: \textbf{y\={a} 6}. Adding twenty coins (\textit{r\={u}pa}), the first person's wealth becomes: \textbf{y\={a} 6, r\={u}pa 20}. The price of ten horses is \textbf{y\={a} 10}. Adding one hundred coins, the second person's wealth becomes: \textbf{y\={a} 10, r\={u}pa 100}.}



\vs



\bilingual{\textbf{समशोधनम्} — एतौ समशोधने छते एव सभी जाते समशोधनार्थं—

\begin{center}\dev{न्यासः---या ६ रूप २०}\\\dev{या १० रूप १००}\end{center}

}{%

\textbf{Equal-clearing} --- These two are made equal for clearing:

\begin{center}\textsc{Setting down:} \textbf{y\={a} 6, r\={u}pa 20}\\\textbf{y\={a} 10, r\={u}pa 100}\end{center}}



\vs



\bilingual{\textbf{नियमः} — अथैकाव्यक्त शोधयेदन्यपक्षादव्ययव्यक्तान्। छोधिते द्वितीयपदमुक्त्वेषु आचयपक्षतुल्यपक्षः शोधितेषु शेषं रूप ४००। अथ कल्पितकलवर्ण \dev{या ४} रूपैः रूप ४०० उक्ते लब्धमेकस्य यावत्तावतो मानं व्यक्तं १००।

}{%

\textbf{Rule} --- One should clear the unknown from one side, and the known from the other side. When cleared, the remainder on the other side is r\={u}pa 400. Then, dividing r\={u}pa 400 by the coefficient of the unknown (y\={a} 4), the value of one y\={a}vatt\={a}vat is obtained as 100. Thus the price of one horse is 100 (coins).}



\newpage



%% ========================================================================

\subsection*{Example: Jewels and Wealth / रत्नधनस्योदाहरणम्}



\bilingual{\begin{sanskritverse}

माणिक्यामल्नीलमौक्तिककिमितिः पञ्चाषसममा—\\

देकस्यान्यतरस्य सप्त नव पट् तद्दलसंख्या सखे।\\

रूपाणां नवतिद्विषष्ठिर्नयोस्तौ तुल्यवित्तौ तथा\\

वीजं प्रतिरुज्ञानि कुरुते मौल्यानि शीघ्रं वद ॥३॥

\end{sanskritverse}

}{%

\begin{flushright}

\textit{Verse 3: Jewels and Equal Wealth}

\end{flushright}

``The measures of ruby, sapphire, pearl, and emerald of one are fifty, seven, nine, and half that number, O friend. Their total wealth is one hundred sixty-two coins each. Tell me quickly the prices, O algebraist, that make them equal.''



\vs

\textsc{Solution:} Let the prices of ruby etc. be y\={a}(1) 3, y\={a} 2, y\={a} 1. The two sides become:

\begin{center}\textbf{y\={a} 19, y\={a} 16, y\={a} 7, r\={u}pa 90}\\\textbf{y\={a} 21, y\={a} 18, y\={a} 6, r\={u}pa 62}\end{center}

By equal-clearing, the value of y\={a}vatt\={a}vat is 4. The prices of jewels are 12, 8, 4.}



\vs



\bilingual{\begin{sanskritverse}

एको द्वीति मम द्विहि शतं धनेन\\

त्वत्तो भवामि हि सखे द्विगुणस्ततोऽस्यः।\\

भूते द्वयोर्यदसि वैनम बहुगुणोऽहं\\

त्वत्तस्तयोर्वद धने मम किं प्रमाणे ॥४॥

\end{sanskritverse}

}{%

\begin{flushright}

\textit{Verse 4: The Merchant's Wealth}

\end{flushright}

``If you give me one hundred, O friend, I become twice as rich as you. But if I give you two hundred, I become some multiple. Tell me, what are our respective measures of wealth?''



\vs

\textsc{Solution:} Let the two merchants' wealth be y\={a} 1 and y\={a} 1. After giving 100, the first becomes twice the second: \textbf{y\={a} 1, r\={u}pa 100} = \textbf{y\={a} 2, r\={u}pa 200}. By clearing: \textbf{y\={a} 4, r\={u}pa 300}, whence y\={a} = 75. The wealths are 40 and 170.}



\newpage



%% ========================================================================

\subsection*{Quadratic and Higher Equations / वर्गसमीकरणम्}



\bilingual{\textbf{उदाहरणम्} — राशिद्वैदशानिशं राशिद्वयनाल्यक्ष कः समो यः स्यात्। राशिकृतिः बहुगुणिता पञ्चत्रिंशद्युता विद्ध ॥६॥

}{%

\textbf{Example} --- ``A certain number when multiplied by twelve and added to its own square equals thirty-five. What is that number?'' Let the number be y\={a} 1. Then: 12y\={a} + y\={a}varga 1 = 35. Clearing gives sides: \textbf{y\={a}v 1, y\={a} 12, r\={u}pa 0} and \textbf{r\={u}pa 35}. Extracting the square root: \textbf{y\={a} 1, r\={u}pa 2} and \textbf{r\={u}pa 3}. From the equation of these, the number is 5.}



\vs



\bilingual{\textbf{वर्गपूरणनियमः} — एषां पूर्वमूलानां च सर्वेषां योगः = \dev{याव ३ या ३१ रूप ८८}। इदमेकाद्वययुतं प्रयोदशवर्गं—

\begin{center}\dev{याव ३ या ३१ रूप ७५}\\\dev{याव ० या ० रूप १६९}\end{center}

}{%

\textbf{Rule for Completion of the Square:} The sum of all these roots = \textbf{y\={a}v 3, y\={a} 31, r\={u}pa 88}. Adding an increment to make a perfect square:

\begin{center}\textbf{y\={a}v 3, y\={a} 31, r\={u}pa 75}\\\textbf{r\={u}pa 169}\end{center}

Equating and multiplying both sides by twelve, taking the double-root from the square 169, the roots are \textbf{y\={a} 6, r\={u}pa 31} and \textbf{r\={u}pa 43}.}



\newpage



%% ========================================================================

\subsection*{The Lotus and the Swallow / हंसपद्मोदाहरणम्}



\bilingual{\textbf{उदाहरणम्} — अत्र वा प्रथमप्रमाणकलं द्वितीयप्रमाणकलं विमले प्रलम्बये तद्गुणगुणितेन द्वितीयमूललेन तुल्यमेव प्रथममूलधनं स्यात्। अतो द्वितीयस्यायं गुणः २।

}{%

\textbf{Example: The Lotus and the Swallow} --- Here, the difference of the two root-measures, multiplied by the product of the multiplier, equals the second root-wealth. The multiplier of the second is 2. If the first root-wealth equals the product of one multiplier and the second root-wealth, divided by the square of the difference, then the difference of measures is 4. Thus the first and second root-wealths are 8 and 4, yielding a result of 2.}



\vs



\bilingual{\begin{sanskritverse}

एकक्रान्तर्चयनात् फलस्य वर्गं विशोध्य परिशिष्यम्।\\

पञ्चक्रान्तेन दत्तं तुल्यः कालः फले च तथैः ॥८॥

\end{sanskritverse}

}{%

\begin{flushright}

\textit{Verse 8: The Lotus Problem}

\end{flushright}

``Having subtracted the square of the result from the first increased [quantity], the remainder, divided by five, gives equal time and result.''



\vs

\textsc{Solution:} The multiplier is 5. Dividing by the decreased multiplier (4), the square of the small difference (16) divides the second wealth, yielding 4. Adding the square of the difference gives the first wealth as 20. By the rule of proportion, time = 20.}



\newpage



%% ========================================================================

\subsection*{Four Merchants with Equal Wealth / चत्वारो वणिजः समधनाः}



\bilingual{\begin{sanskritverse}

माणिक्याणुकणिनीलादीनां मुक्ताफलानां शतं\\

षड्भागि च पञ्च रत्नवीरां येषां चतुर्णां धनम्।\\

संग्रहे हयरैन ते निजधनाद्वेकैक मिथो\\

जातास्तुल्यधनाः पृथग्वद सखे तद्दनमैल्यानि मे ॥२॥

\end{sanskritverse}

}{%

\begin{flushright}

\textit{Verse 2: Four Merchants of Equal Wealth}

\end{flushright}

``Of rubies, coral, sapphires, pearls, and emeralds, there are one hundred; six parts and five jewels. Four merchants whose wealth, when one gives to another, they become equal. Tell me separately, O friend, the prices of their wealth.''



\vs

\textsc{Method:} The unknowns are postulated as y\={a}vatt\={a}vats. After four equal-clearings, the root is obtained by intelligent reasoning. As Bh\={a}skara himself says:

\begin{quote}

``Neither squared nor unsquared is the seed; seeds are not separate. Intelligence alone is the seed, for endless are the suppositions.''

\end{quote}}



\newpage



%% ========================================================================

\subsection*{Geometric Application: Equilateral Triangle / समभुजत्रिकोणस्य क्षेत्रव्यवहारः}



\bilingual{\textbf{न्यासः} — क्षेत्रव्यवहारात् समभुजत्रिकोणस्य जाते लम्बवर्गः = \dev{याव १ रूप १५, रूप २००}। या १ या १ क १८ रूप १। द्वितीयाव्यावावर्गी = \dev{याव १ या क ६२ या २ रूप १९ क ६२}।

}{%

\textbf{Setting Down:} From geometric computation for an equilateral triangle, the square of the altitude is: \textbf{y\={a}v 1, r\={u}pa 15, r\={u}pa 200}, or \textbf{y\={a} 1, k\={a} 18, r\={u}pa 1}. The second altitude-square is: \textbf{y\={a}v 1, y\={a} k\={a} 62, y\={a} 2, r\={u}pa 19, k\={a} 62}. Subtracting the base 6, the second altitude-square becomes: \textbf{y\={a}v 1, y\={a} 2, y\={a} k\={a} 62, r\={u}pa 13, k\={a} 62}.}



\vs



\bilingual{\textbf{समशोधनम्} — एतौ समाविति समशोधने छते जाते पक्षौ—

\begin{center}\dev{र २५ क ५१२}\\\dev{या २ या क ६२}\end{center}

}{%

\textbf{Equal-clearing:} These being made equal, the sides after clearing are:

\begin{center}\textbf{r\={u}pa 25, k\={a} 512}\\\textbf{y\={a} 2, y\={a} k\={a} 62}\end{center}}



\newpage



%% ========================================================================

\section*{Equations with Unknown Squares / अव्यक्तवर्गादिसमीकरणम्}

\addcontentsline{toc}{section}{Avyakta-varg\={a}di-sam\={\i}kara\d{n}a (pp.~77--95)}



\bilingual{\textbf{सूत्रम्} —

\begin{sanskritverse}

अव्यक्तवर्गादि यदा द्वयोः पक्षौ तदैकेन निःशेषं किंचित्।\\

क्षेपं तयोरेन पदप्रदः स्याद्व्यक्तः सहेत परेन भूयः ॥१॥

\end{sanskritverse}

\begin{sanskritverse}

व्यक्तस्य मूलस्य समीकृते समव्यक्तमानं बहु लभ्यते तत्।\\

न निर्वहेत् चेद्धनवर्गयोस्त्वेति तदा क्षेपमिदं स्यबुद्ध्या ॥२॥

\end{sanskritverse}

}{%

\textbf{Rule (Verses 1--2):} ---

\begin{flushright}

\textit{When unknown squares appear on both sides,}\\

\textit{divide by something to make one side a perfect square.}\\

\textit{The additive must be such that it yields}\\

\textit{a rational root when combined with the known term.}\\[0.5em]

\textit{When the root of the known is equated,}\\

\textit{many unknown values may be obtained;}\\

\textit{if the square does not succeed,}\\

\textit{then by intelligence determine the additive.}

\end{flushright}}



\vs



\bilingual{\begin{sanskritverse}

अव्यक्तमूलर्कांक ततोऽस्य व्यक्तस्य पक्षस्य पदं यदि स्यात्।\\

मूलं धनं तथा विधाय साव्यमव्यक्तमानं द्विधा कल्पितं स्यात्॥३॥

\end{sanskritverse}

}{%

\textbf{Verse 3:}

\begin{flushright}

\textit{If the root of the unknown side, having a number,}\\

\textit{when made the root of the known side---}\\

\textit{having made that root positive or negative,}\\

\textit{the unknown value is to be assumed twofold.}

\end{flushright}}



\newpage



%% ========================================================================

\subsection*{Example: Series of Numbers / राशिश्रेण्युदाहरणम्}



\bilingual{\begin{sanskritverse}

स्वार्धपञ्चाशनवमैर्युक्ताः के स्युः समाश्रयाः।\\

अन्यैर्द्वादशहीनाख्य परिशिष्टाख्य तान् वद ॥१४॥

\end{sanskritverse}

}{%

\begin{flushright}

\textit{Verse 14: Numbers with Fractional Relations}

\end{flushright}

``What numbers, when increased by their own half, five-ninths, and [other fractions], become equal? And when diminished by others, twelve less---tell me the remainders.''



\vs

\textsc{Solution:} Let the equal number be y\={a}vatt\={a}vat 1. By the reverse process, the numbers are y\={a} $\frac{2}{5}$, y\={a} $\frac{5}{6}$, y\={a} $\frac{9}{10}$. Dividing by the greatest common measure and multiplying by y\={a}vatt\={a}vat value 150, the numbers are 100, 125, 135.}



\vs



\bilingual{\begin{sanskritverse}

त्रयोदश तथा पञ्च करण्यो मूजयोनितो।\\

भूजाता च चत्वारः फले भूमिं वदाऽखु मे ॥१५॥

\end{sanskritverse}

}{%

\begin{flushright}

\textit{Verse 15: The Excavated Plot}

\end{flushright}

``The square-roots of thirteen and five, diminished by the base, and four as the result---tell me the area, O algebraist.''



\vs

\textsc{Solution:} Let the area be y\={a}vatt\={a}vat 1. By the formula ``altitude times half-base divided [equals the area of a triangle],'' the section-ratio is k\={a} $\frac{5}{13}$. Its square is k\={a} 5. Subtracting from r\={u}pa 5, r\={u}pa $\frac{5}{13}$, the root gives the altitude k\={a} $\frac{5}{13}$.}



\newpage



%% ========================================================================

\subsection*{Geometric Diagrams / क्षेत्रचित्राणि}



\bilingual{\textbf{क्षेत्राणि} — तेषां स्वरूपाणि यथा—

\begin{center}

\begin{tikzpicture}[scale=0.5]

\draw[thick] (0,0) grid (6,6);

\draw[very thick] (0,0) rectangle (6,6);

\node at (3,-0.5) {\dev{६}};

\node at (-0.5,3) {\dev{६}};

\node at (6.5,3) {\dev{६}};

\node at (3,6.5) {\dev{६}};

\end{tikzpicture}

\hspace{0.5cm}

\begin{tikzpicture}[scale=0.5]

\draw[thick] (0,0) grid (3,3);

\draw[very thick] (0,0) rectangle (3,3);

\node at (1.5,-0.5) {\dev{३}};

\node at (-0.5,1.5) {\dev{३}};

\node at (3.5,1.5) {\dev{३}};

\node at (1.5,3.5) {\dev{३}};

\end{tikzpicture}

\end{center}

}{%

\textbf{Geometric Figures:} Their forms are as follows:

\begin{center}

\textit{Square of side 6}\hspace{0.5cm}\textit{Square of side 3}

\end{center}

\textit{The large square is 36 units in area; the small square is 9 units. The difference is 27, which may be analyzed as the sum of gnomon regions.}



\vs



\bilingual{\begin{sanskritverse}

चतुर्गुणस्य घातस्य युतिवर्गस्य चान्तरम्।\\

राश्यन्तरकृतेस्तुल्यं द्वयोरन्त्यकयोर्यथा ॥१७॥

\end{sanskritverse}

}{%

\textbf{Verse 17:}

\begin{flushright}

\textit{The difference between the square of the sum}\\

\textit{and four times the product}\\

\textit{equals the square of the difference of the two numbers,}\\

\textit{as between the last two terms.}

\end{flushright}

\textsc{Numerical verification:} For numbers 3 and 5: $(3+5)^2 - 4 \cdot 3 \cdot 5 = 64 - 60 = 4 = (5-3)^2$. This is the algebraic identity: $(a+b)^2 - 4ab = (a-b)^2$.}



\newpage



%% ========================================================================

\part*{Chapter II: Equations with Many Unknowns}

\addcontentsline{toc}{part}{Chapter II: Equations with Many Unknowns}



\section*{\dev{अनेकवर्णसमीकरणम्} $\parallel$ Anekavar\d{n}a-sam\={\i}kara\d{n}a}

\addcontentsline{toc}{section}{Anekavar\d{n}a-sam\={\i}kara\d{n}a (pp.~96--114)}



\bilingual{\lettrine[lines=2, lraise=0.1, nindent=0em, slope=-.5em]{\dev{अ}}{}थानेकवर्णसमीकरणम्। अत्र बहूनां यावत्तावतां सहचरितानां समीकरणं दीयते। यदा बहूनामज्ञातानां राशीनां मूल्यानि निर्णेतव्यानि भवन्ति, तदा प्रत्येकं राश्यै यावत्तावत् प्रकल्प्यते। तदनन्तरं कुट्टकविधिना मानानि निर्णीयन्ते।

}{%

\lettrine[lines=2, lraise=0.1, nindent=0em, slope=-.5em]{N}{}ow, equations with many unknowns. Here, equations are given involving several associated y\={a}vatt\={a}vats. When the values of many unknown quantities must be determined, a y\={a}vatt\={a}vat is postulated for each quantity. Thereafter, the values are determined by the \textit{ku\d{t}\d{t}aka} (pulverizer) method.}



\vs



\bilingual{\textbf{मुख्यनियमः} — अत्र माणिक्यादीनां मौल्यानि यावत्तावदादीनि प्रकल्प्य तद्गुणपरितसंख्यां च छत्वा रूपाणि च प्रक्षिप्य समशोधनार्थं—

\begin{center}

\dev{न्यासः---या ५ का ८ नी ३ रूप १०}\\

\dev{या ३ का ९ नी ५ रूप ६२}

\end{center}

}{%

\textbf{Main Rule:} Let the prices of rubies etc. be postulated as y\={a}vatt\={a}vats etc., and multiplying by their respective coefficients and adding r\={u}pas, for equal-clearing:

\begin{center}

\textsc{Setting down:} \textbf{y\={a} 5, k\={a} 8, n\={\i} 3, r\={u}pa 10}\\

\textbf{y\={a} 3, k\={a} 9, n\={\i} 5, r\={u}pa 62}

\end{center}}



\vs



\bilingual{\textbf{समशोधनम्} — अत्र वर्गं शोध्येदित्यादिना जाता यावत्तावद्विमतिः:

\begin{center}\dev{या $= \dfrac{\text{का\textsuperscript{`}}~\text{नी}~\text{१ र}~\text{२८}}{\text{२}}$}\end{center}

}{%

\textbf{Equal-clearing:} Clearing the square etc., the value of y\={a}vatt\={a}vat is:

\begin{center}$\displaystyle\text{y\={a}} = \frac{-\text{k\={a}} + \text{n\={\i}} + \text{r\={u}pa }28}{2}$\end{center}}



\newpage



%% ========================================================================

\subsection*{Tabular Solution / सारणीकृतसाधनम्}



\bilingual{\textbf{निर्माणप्रक्रिया} — प्रथमं नी=३ कल्पिते यदि तत्र पीतकमानं किमपि त्रिगुणितमेव। तदा—

\begin{center}

\begin{tabular}{rclcrclcrcl}

\dev{यदि} & \dev{पी=३९} & \dev{,} & \dev{का=४५} & \dev{अतः} & \dev{या=१} & \dev{(१)} \\

\dev{,} & \dev{पी=४२} & \dev{,} & \dev{का=३०} & \dev{,} & \dev{या=२१} & \dev{(२)} \\

\dev{,} & \dev{पी=४५} & \dev{,} & \dev{का=१५} & \dev{,} & \dev{या=३३} & \dev{(३)} \\

\end{tabular}

\end{center}

}{%

\textbf{Construction Procedure:} First, if n\={\i}=3 is postulated, the yellow value must be some multiple of three. Then:

\begin{center}

\begin{tabular}{rclcrclcrcl}

If & p\={\i}=39 & $\Rightarrow$ & k\={a}=45 & then & y\={a}=1 & (1)\\

If & p\={\i}=42 & $\Rightarrow$ & k\={a}=30 & then & y\={a}=21 & (2)\\

If & p\={\i}=45 & $\Rightarrow$ & k\={a}=15 & then & y\={a}=33 & (3)\\

\end{tabular}

\end{center}

\textsc{Note:} This tabular method generates an infinite family of solutions. The values are obtained by the \textit{ku\d{t}\d{t}aka} (pulverizer) algorithm applied to the reduced equation.}



\newpage



%% ========================================================================

\subsection*{The General Rule / सामान्यसूत्रम्}



\bilingual{\begin{sanskritverse}

द्वयोः समीकरणयोर्मूलमेकं यदि विद्यते।\\

तदा तदन्यत्समीकर्य मूलं तद्वदनुष्ठितम्॥

\end{sanskritverse}

}{%

\textbf{General Rule:}

\begin{flushright}

\textit{When the root of one of two equations is known,}\\

\textit{then having equated the other,}\\

\textit{its root is obtained by the same procedure.}

\end{flushright}

\textsc{Explanation:} This is the fundamental principle of elimination. If one variable can be expressed from one equation, substitute this expression into the second equation. This is the \textit{madhyam\={a}hara\d{n}a} (elimination of the middle) method in its general form.}



\newpage



%% ========================================================================

\section*{\dev{अनेकवर्णमध्यमाहरणभेदाः} $\parallel$ Anekavar\d{n}a-madhyam\={a}hara\d{n}a-bhed\={a}h}

\addcontentsline{toc}{section}{Anekavar\d{n}a-madhyam\={a}hara\d{n}a-bhed\={a}h (pp.~107--114)}



\bilingual{\lettrine[lines=2, lraise=0.1, nindent=0em, slope=-.5em]{\dev{अ}}{}थानेकवर्णमध्यमाहरणभेदाः। अत्र बहुवर्णसमीकरणानां मध्यमाहरणस्य विभिन्नाः पद्धतयः प्रदर्श्यन्ते। मध्यमाहरणमित्युक्ते अनेकवर्णसमीकरणेभ्यो मध्यमवर्णस्य निरसनं, अन्येषां च मानानां ततो निर्णयः।

}{%

\lettrine[lines=2, lraise=0.1, nindent=0em, slope=-.5em]{N}{}ow, variations of the elimination method for equations with many unknowns. Here, various methods of \textit{madhyam\={a}hara\d{n}a} (elimination of the middle [term]) for equations with many unknowns are demonstrated. By \textit{madhyam\={a}hara\d{n}a} is meant the elimination of the middle unknown from a system of equations, and the determination of the remaining quantities thereafter.}



\vs



\bilingual{\begin{sanskritverse}

इष्टं तयोः प्रथमपक्षादेन\\

छत्तोकवत् प्रथमवर्णमितिकृत् साध्या।\\

हृद् यं भवेत् प्रकृतिवर्गमितिः सुखोर्मि-\\

र्वै कृतिप्रकृतित्रय नियोजनीया ॥५॥

\end{sanskritverse}

}{%

\textbf{Verse 5:}

\begin{flushright}

\textit{Having chosen [a number] from the first side of the two,}\\

\textit{it should be made to have the measure of the first unknown,}\\

\textit{as if divided. That which would be in the heart}\\

\textit{[i.e., the divisor]---the square-root of the original,}\\

\textit{three [terms] involving the square and the original}\\

\textit{must be applied.}

\end{flushright}

\textsc{Explanation:} This verse describes a sophisticated elimination procedure. The ``three terms involving the square and the original'' refers to the quadratic, linear, and constant terms that arise when completing the square during elimination.}



\newpage



%% ========================================================================

\subsection*{Example: Doubled Number with Square / द्विगुणितराशिवर्गयुतः}



\bilingual{\begin{sanskritverse}

को राशिद्विगुणो राशिवर्गेः पट्डिमः समन्वितः।\\

मूलद्वो जायते बीजगणितं वद तत् त्वम् ॥१॥

\end{sanskritverse}

}{%

\begin{flushright}

\textit{Verse 1: A Number with Its Square}

\end{flushright}

``What number, when doubled and added to its square-root's square (i.e., the number itself), produces two roots? Speak that algebra, O you.''



\vs

\textsc{Solution:} Let the first unknown's doubled quantity with its square be y\={a}v 6, y\={a} 2. For the square, the setting-down is:

\begin{center}\textbf{y\={a}v 6, y\={a} 2, k\={a}v 0}\\\textbf{y\={a}v 0, y\={a} 0, k\={a}v 1}\end{center}

By equal-clearing, the sides become y\={a}v 6, y\={a} 2, k\={a}v 1. Multiplying by the coefficient and adding r\={u}pa, the first root is \textbf{y\={a} 6, r\={u}pa 1}. The second root is \textbf{k\={a} 2}, or k\={a} 20 with residue 49. Equating with the first side y\={a} 6, r\={u}pa 1, the value of y\={a}vatt\={a}vat is $\frac{2}{5}$ or 8.}



\newpage



%% ========================================================================

\subsection*{Example: Sum and Product / राशियोगघातोदाहरणम्}



\bilingual{\begin{sanskritverse}

राशियोगकृतिर्मित्रा राश्योर्गुणघनेन चेत्।\\

द्विस्य धनयोगस्य सा तुल्या गणकौच्यताम् ॥२॥

\end{sanskritverse}

}{%

\begin{flushright}

\textit{Verse 2: The Square of the Sum}

\end{flushright}

``The square of the sum of numbers, O friend, when equal to the cube of their product---let the calculator declare the sum of their doubled wealth.''



\vs

\textsc{Solution:} The procedure by which [this is done]---intelligent ones should postulate numbers. Postulated: (y\={a} 1, k\={a} 1), (y\={a} 1, k\={a} 1). Their sum is 2. By its double, the wealth is y\={a}v 8, y\={a}v 4. The individual wealths of the numbers: y\={a}v 1, y\={a} k\={a} m\={a} 2, k\={a}v y\={a} m\={a} 2, k\={a}v m\={a} 1. For the second: bh\={a}ga 1, y\={a} k\={a} m\={a} 2, k\={a}v y\={a} m\={a} 2, k\={a}v m\={a} 1. Their sum is y\={a}v 2, k\={a}v y\={a} m\={a} 6. Twice this: y\={a}v 4, k\={a}v y\={a} m\={a} 12, for equal-clearing.}



\newpage



%% ========================================================================

\subsection*{The Elimination Procedure / मध्यमाहरणप्रक्रिया}



\bilingual{\textbf{मुख्यन्यासः} — अत्र माणिक्यादीनां मौल्यानि यावत्तावदादीनि प्रकल्प्य तद्गुणपरितसंख्यां च छत्वा रूपाणि च प्रक्षिप्य समशोधनार्थं—

\begin{center}

\dev{न्यासः---या ५ का ८ नी ३ रूप १०}\\

\dev{या ३ का ९ नी ५ रूप ६२}

\end{center}

}{%

\textbf{Main Setting-down:} Postulating the prices of rubies etc. as y\={a}vatt\={a}vats, multiplying by their respective coefficients and adding r\={u}pas, for equal-clearing:

\begin{center}

\textsc{Setting:} \textbf{y\={a} 5, k\={a} 8, n\={\i} 3, r\={u}pa 10}\\

\textbf{y\={a} 3, k\={a} 9, n\={\i} 5, r\={u}pa 62}

\end{center}}



\vs



\bilingual{\textbf{समशोधनम्} — अत्र वर्गं शोध्येदित्यादिना जाता यावत्तावद्विमतिः:

\begin{center}\dev{या $= \dfrac{\text{का\textsuperscript{`}}~\text{नी}~\text{१ र}~\text{२८}}{\text{२}}$}\end{center}

}{%

\textbf{Clearing:} Clearing the squares, the value of y\={a}vatt\={a}vat is:

\begin{center}$\displaystyle\text{y\={a}} = \frac{-\text{k\={a}} + \text{n\={\i}} + \text{r\={u}pa }28}{2}$\end{center}

\textsc{Note:} This is a single equation with three unknowns. By the \textit{ku\d{t}\d{t}aka} method, infinitely many solutions are obtained.}



\newpage



%% ========================================================================

\subsection*{Concluding Verses / समाप्तिश्लोकाः}



\bilingual{\begin{sanskritverse}

नवमिः सममिः शुण्णः को राशिकिंवशात् हतः।\\

यद्वैक्यं फले काल्यं भवेत् बहुविशतेऽमितम् ॥१०॥

\end{sanskritverse}

}{%

\begin{flushright}

\textit{Verse 10: Equal Remainder}

\end{flushright}

``Nine is the equal remainder. What number, when multiplied by a certain divisor, when added to another and divided by the divisor, yields an equal remainder of twenty?''



\vs

\textsc{Solution:} The divisor is postulated as r\={u}pa 16. The number is l\={a} 1. Having obtained one measure by the single-divisor reduction, the multiplier becomes the product: y\={a} 16, k\={a} 30. This result, diminished by the divisor: y\={a} 16, k\={a} 26, equated to twenty, by \textit{ku\d{t}\d{t}aka} the value y\={a}vatt\={a}vat is n\={\i} 29, r\={u}pa 297. Here, since only one sum of results is indicated, no additive is required.}



\vs



\bilingual{\begin{sanskritverse}

कलितसनवशुण्णो राशिकिंवशद्विभाजितः।\\

यद्वैक्यमपि त्रिशद्दत्तमेकादशाऽऽक्रमम् ॥११॥

\end{sanskritverse}

}{%

\begin{flushright}

\textit{Verse 11: Divided with Remainder}

\end{flushright}

``A number increased by nine and diminished by something, divided by a certain divisor, when added together gives thirty as the result in the eleventh step.''



\vs

\textsc{Solution:} Let the first number be y\={a}vatt\={a}vat plus the known. The second is 6y\={a} + 2. By the statement of the first, the value of y\={a}vatt\={a}vat = $\frac{3\text{k\={a}} + \text{known}}{6}$, and by the third statement: y\={a}} = $\frac{9\text{n\={\i}} + 3 - \text{known}}{6}$. By the \textit{ku\d{t}\d{t}aka} method, the first number is postulated as 36, the second as 126, r\={u}pa 104.}



\newpage



%% ========================================================================

\section*{The Cakrav\={a}la (Cyclic) Method $\parallel$ चक्रवालपद्धतिः}

\addcontentsline{toc}{section}{The Cakrav\={a}la Method (Vargaprak\d{r}ti)}



\bilingual{\textbf{चक्रवालपद्धतिः} — \dev{वर्गप्रकृति} इत्युक्ते $x^2 - Ny^2 = 1$ रूपं समीकरणं यत्र $N$ अज्ञातवर्गे कारः (गुणकः), $x$ प्रकृतिः (ज्येष्ठमूलम्), $y$ ह्रस्वम् (कनिष्ठमूलम्)। भास्कराचार्यस्य \textbf{चक्रवालपद्धतिः} अस्मिन् समीकरणे साधने सर्वोत्तमा अल्गोरिद्मिकपद्धतिः अस्ति। इयं पद्धतिः सदैव परिसमाप्तिं गच्छति तथा च न्यूनतमं समाधानं प्रयच्छति।

}{%

\textbf{The Cakrav\={a}la (Cyclic) Method:} The \textit{vargaprak\d{r}ti} is the equation $x^2 - Ny^2 = 1$, where $N$ is the \textit{k\={a}ra\d{n}a} (multiplier/prak\d{r}ti), $x$ is the \textit{jye\d{s}\d{t}ha-m\={u}la} (greater root), and $y$ is the \textit{kani\d{s}\d{t}ha-m\={u}la} (lesser root). Bh\={a}skara's \textbf{Cakrav\={a}la} (``cyclic'') method is the supreme algorithmic procedure for solving this equation. This method always terminates and produces the minimal solution.}



\vs



\bilingual{\textbf{चक्रवालस्य पद्धतिः} —

\begin{enumerate}[label=\arabic*.]

\item मूलस्य ($\sqrt{N}$) समीपतमं पूर्णाङ्कं गृहाण।

\item प्रथमं समीकरणं रचय: $a^2 - N \cdot 1^2 = k$, यत्र $a$ समीपतमपूर्णाङ्कं, $k$ क्षेपः।

\item ततः $b = \dfrac{a + m}{|k|}$ गणय, यत्र $m$ अनेन प्रकारेण वृत्यते यत् $a + m \equiv 0 \pmod{|k|}$ तथा च $m^2 - N$ न्यूनतमं भवति।

\item नवीन समीकरणं रचय: $a_{\text{new}}^2 - N \cdot b^2 = k_{\text{new}}$।

\item यदा $k_{\text{new}} = 1, 2, -1, -2, 4, -4$ भवति, तदा \textit{भावना} (composition) क्रियते।

\item अन्यथा चक्रं पुनः आरभ्यते।

\end{enumerate}

}{%

\textbf{The Cakrav\={a}la Procedure:}

\begin{enumerate}[label=\arabic*.]

\item Take the integer nearest to $\sqrt{N}$.

\item Form the first equation: $a^2 - N \cdot 1^2 = k$, where $a$ is the nearest integer, $k$ is the additive (\textit{k\d{s}epa}).

\item Compute $b = \dfrac{a + m}{|k|}$, where $m$ is chosen such that $a + m \equiv 0 \pmod{|k|}$ and $m^2 - N$ is minimized.

\item Form the new equation: $a_{\text{new}}^2 - N \cdot b^2 = k_{\text{new}}$.

\item When $k_{\text{new}} = 1, 2, -1, -2, 4, -4$, apply \textit{bh\={a}van\={a}} (composition).

\item Otherwise, repeat the cycle.

\end{enumerate}}



\newpage



%% ========================================================================

\subsection*{The Famous Example $N=61$ / प्रसिद्धमुदाहरणम् $N=61$}



\bilingual{\textbf{उदाहरणम्} — \dev{``एकषष्टिप्रकृतिः''} इति प्रसिद्धमुदाहरणं भास्कराचार्यस्य अस्ति। अत्र $N = 61$। एतत् समीकरणं $x^2 - 61y^2 = 1$ इत्यस्य न्यूनतमं समाधानं अत्यन्तं विस्तृतं जातं यत् पूर्वं यावद्ब्रह्मगुप्तस्य कालपर्यन्तं कठिनतमं मन्यते स्म।

}{%

\textbf{Example:} The ``\textit{prak\d{r}ti} sixty-one'' is Bh\={a}skara's most celebrated example. Here $N = 61$. The minimal solution of $x^2 - 61y^2 = 1$ is enormously large, and was considered the most difficult case from the time of Brahmagupta until Bh\={a}skara.}



\vs



\bilingual{\textbf{चक्रवालस्य परिक्रमाः} —

\begin{center}

\begin{tabular}{|c|c|c|c|c|c|}

\hline

\textbf{\dev{चक्रम्}} & \textbf{\dev{गुणकः}} & \textbf{\dev{ज्येष्ठम्}} & \textbf{\dev{कनिष्ठम्}} & \textbf{\dev{क्षेपः}} & \textbf{\dev{वृत्तिः}} \\

\hline

\dev{१} & \dev{६१} & \dev{८} & \dev{१} & \dev{-१३} & \dev{प्रारम्भः} \\

\dev{२} & \dev{६१} & \dev{४७} & \dev{६} & \dev{७} & \dev{चक्रम्} \\

\dev{३} & \dev{६१} & \dev{४८} & \dev{७} & \dev{-१३} & \dev{चक्रम्} \\

\dev{४} & \dev{६१} & \dev{६३} & \dev{८} & \dev{-४} & \dev{भावना} \\

\dev{५} & \dev{६१} & \dev{११९} & \dev{१५} & \dev{४} & \dev{भावना} \\

\dev{६} & \dev{६१} & \dev{३९७} & \dev{५१} & \dev{-४} & \dev{भावना} \\

\dev{७} & \dev{६१} & \dev{१३२३} & \dev{१७०} & \dev{७} & \dev{चक्रम्} \\

\dev{८} & \dev{६१} & \dev{१३२४} & \dev{१७१} & \dev{-१३} & \dev{चक्रम्} \\

\dev{९} & \dev{६१} & \dev{१८९१} & \dev{२४२} & \dev{७} & \dev{चक्रम्} \\

\dev{१०} & \dev{६१} & \dev{२८५६} & \dev{३६५} & \dev{-३} & \dev{चक्रम्} \\

\dev{११} & \dev{६१} & \dev{२८५७} & \dev{३६६} & \dev{-१} & \dev{समाप्तिः} \\

\hline

\end{tabular}

\end{center}

}{%

\textbf{Cycles of the Cakrav\={a}la:}

\begin{center}

\begin{tabular}{|c|c|c|c|c|c|}

\hline

\textbf{Cycle} & \textbf{Prak\d{r}ti $N$} & \textbf{Jye\d{s}\d{t}ha $a$} & \textbf{Kani\d{s}\d{t}ha $b$} & \textbf{K\d{s}epa $k$} & \textbf{Operation} \\

\hline

1 & 61 & 8 & 1 & $-13$ & Initiation \\

2 & 61 & 47 & 6 & 7 & Cycle \\

3 & 61 & 48 & 7 & $-13$ & Cycle \\

4 & 61 & 63 & 8 & $-4$ & Bh\={a}van\={a} \\

5 & 61 & 119 & 15 & 4 & Bh\={a}van\={a} \\

6 & 61 & 397 & 51 & $-4$ & Bh\={a}van\={a} \\

7 & 61 & 1323 & 170 & 7 & Cycle \\

8 & 61 & 1324 & 171 & $-13$ & Cycle \\

9 & 61 & 1891 & 242 & 7 & Cycle \\

10 & 61 & 2856 & 365 & $-3$ & Cycle \\

11 & 61 & 2857 & 366 & $-1$ & \textbf{Termination} \\

\hline

\end{tabular}

\end{center}}



\newpage



%% ========================================================================

\subsection*{Final Composition: The Solution / अन्तिमभावना: समाधानम्}



\bilingual{\textbf{समाधानम्} — यदा क्षेपः $k = -1$ प्राप्यते (चक्रे ११), तदा \textit{भावनया} सह सम्पूर्णं समाधानं प्राप्यते। अन्तिमे समीकरणे $2857^2 - 61 \cdot 366^2 = -1$ भवति। अस्य वर्गीकरणेन ($k = -1$ इत्यस्य द्विवारं प्रयोगेन) $k = 1$ प्राप्यते।



\vspace{0.5em}

\textbf{\dev{अतः न्यूनतमं समाधानम्:}}

\begin{center}

\fbox{\dev{ज्येष्ठमूलम् $x = 1{,}766{,}319{,}049$}}\\[0.3em]

\fbox{\dev{कनिष्ठमूलम् $y = 226{,}153{,}980$}}

\end{center}

\vspace{0.3em}

\textbf{\dev{सत्यापनम्:}}

\begin{center}

$1{,}766{,}319{,}049^2 - 61 \cdot 226{,}153{,}980^2 = 1$

\end{center}

}{%

\textbf{The Solution:} When the additive $k = -1$ is reached (at cycle 11), the complete solution is obtained by \textit{bh\={a}van\={a}} (composition). The final equation yields $2857^2 - 61 \cdot 366^2 = -1$. Squaring this (applying $k = -1$ twice) gives $k = 1$.



\vspace{0.5em}

\textbf{The minimal solution is:}

\begin{center}

\fbox{Greater root $x = 1{,}766{,}319{,}049$}\\[0.3em]

\fbox{Lesser root $y = 226{,}153{,}980$}

\end{center}

\vspace{0.3em}

\textbf{Verification:}

\begin{center}

$1{,}766{,}319{,}049^2 - 61 \cdot 226{,}153{,}980^2 = 1$ $\checkmark$

\end{center}



\vs

\textsc{Historical Note:} This solution was not matched in Europe until the work of Fermat, Brouncker, and Wallis in the 17th century. Bh\={a}skara's Cakrav\={a}la algorithm always terminates---a fact that was only rigorously proved in the 20th century by Krishnaswami Alladi. The method is superior to the European ``chain method'' of Brouncker in both elegance and efficiency.}



\newpage



%% ========================================================================

\section*{The Bh\={a}vita Method $\parallel$ भावितपद्धतिः}

\addcontentsline{toc}{section}{The Bh\={a}vita Method (Product Equations)}



\bilingual{\textbf{भावितपद्धतिः} — \dev{भावित} इत्युक्ते द्वयोः अज्ञातराश्योः घातः (product)। यदा समीकरणे द्वयोः वर्णयोः गुणनफलं प्रविष्टं भवति, तदा \textit{भावितपद्धतिः} प्रयुज्यते। अस्याः पद्धतेर्मुख्यं प्रयोजनं \textit{वर्गप्रकृतेः} समाधानेषु \textit{भावनायाः} (composition) क्रियायां विद्यते।

}{%

\textbf{The Bh\={a}vita Method:} \textit{Bh\={a}vita} means the product of two unknown quantities. When the product of two unknowns appears in an equation, the \textit{bh\={a}vita} method is applied. The principal application of this method is in the \textit{bh\={a}van\={a}} (composition) operation for solving the \textit{vargaprak\d{r}ti}.}



\vs



\bilingual{\textbf{भावनानियमः} — यदि $(a, b, k)$ तथा $(A, B, K)$ द्वे वर्गप्रकृतिसमाधाने स्तः, अर्थात् $a^2 - Nb^2 = k$ तथा $A^2 - NB^2 = K$, तदा भावनया नवीनं समाधानं प्राप्यते:

\begin{center}

\fbox{$x = aA \pm NbB$, \quad $y = aB \pm Ab$}

\end{center}

\vspace{0.3em}

यस्य क्षेपः नवीनं समीकरणं $x^2 - Ny^2 = kK$ सन्तोषयति।

}{%

\textbf{The Bh\={a}van\={a} Rule:} If $(a, b, k)$ and $(A, B, K)$ are two solutions of the vargaprak\d{r}ti, i.e., $a^2 - Nb^2 = k$ and $A^2 - NB^2 = K$, then by \textit{bh\={a}van\={a}} a new solution is obtained:

\begin{center}

\fbox{$x = aA \pm NbB$, \quad $y = aB \pm Ab$}

\end{center}

\vspace{0.3em}

whose additive satisfies the new equation $x^2 - Ny^2 = kK$.



\vs

\textsc{Special Case:} When $k = K = 1$, the product $kK = 1$, so the composed solution directly solves $x^2 - Ny^2 = 1$. When $k = K = -1$, then $kK = 1$ again, so two solutions with additive $-1$ compose to a solution with additive $+1$. This is the key to terminating the Cakrav\={a}la method.}



\newpage



%% ========================================================================

\section*{Summary of Part II / द्वितीयभागस्य संक्षेपः}

\addcontentsline{toc}{section}{Summary}



\bilingual{\textbf{समीकरणविधयः} — अस्य भागे निम्नलिखिताः प्रधानगणितविधयः उपदिष्टाः सन्ति:



\begin{enumerate}

\item \textbf{एकवर्णसमीकरणम्} — रेखीयसमीकरणानां, वर्गसमीकरणानां च साधनम्।

\item \textbf{अव्यक्तवर्गादिसमीकरणम्} — अज्ञातवर्गयुक्तसमीकरणानां साधनम्।

\item \textbf{अनेकवर्णसमीकरणम्} — बहुवर्णसमीकरणानां कुट्टकविधिना साधनम्।

\item \textbf{मध्यमाहरणम्} — मध्यमवर्णस्य निरसनपद्धतिः।

\item \textbf{चक्रवालपद्धतिः} — वर्गप्रकृतेः ($x^2 - Ny^2 = 1$) सर्वोत्तमा साधनपद्धतिः।

\item \textbf{भावितपद्धतिः} — द्विवर्णघातसमीकरणानां तथा भावनायाः क्रिया।

\end{enumerate}

}{%

\textbf{Equation Methods:} In this part, the following principal mathematical methods are taught:



\begin{enumerate}

\item \textbf{Equations with One Unknown} --- solution of linear and quadratic equations.

\item \textbf{Equations with Unknown Squares} --- solution of equations containing unknown squares.

\item \textbf{Equations with Many Unknowns} --- solution by the \textit{ku\d{t}\d{t}aka} (pulverizer) method.

\item \textbf{Elimination of the Middle} --- the method of eliminating the middle unknown.

\item \textbf{The Cakrav\={a}la Method} --- the supreme method for solving $x^2 - Ny^2 = 1$.

\item \textbf{The Bh\={a}vita Method} --- product equations and the composition operation.

\end{enumerate}}



\vs



\bilingual{\textbf{भास्करस्य योगदानम्} — भास्कराचार्यस्य \textbf{चक्रवालपद्धतिः} गणितस्य इतिहासे अत्यन्तं महत्त्वपूर्णा अस्ति। इयं पद्धतिः—

\begin{itemize}

\item सदैव परिसमाप्तिं गच्छति (always terminates)

\item न्यूनतमं समाधानं प्रयच्छति (produces the minimal solution)

\item युरोपीयपद्धतेः अपेक्षया अधिकतया कुशला (more efficient than European methods)

\item आधुनिकअल्गोरिद्मिकविश्लेषणस्य आधारभूता (foundational for modern algorithmic analysis)

\end{itemize}

}{%

\textbf{Bh\={a}skara's Contribution:} Bh\={a}skara's \textbf{Cakrav\={a}la method} is one of the most significant achievements in the history of mathematics. This method:

\begin{itemize}

\item Always terminates (always terminates)

\item Produces the minimal solution (produces the minimal solution)

\item Is more efficient than European methods of the 17th century (more efficient than European methods)

\item Is foundational for modern algorithmic analysis (foundational for modern algorithmic analysis)

\end{itemize}}



\newpage



%% ========================================================================

\vspace*{2cm}



\begin{center}

{\Large\bfseries\color{sectioncolor} End of Part II $\parallel$ द्वितीयभागसमाप्तिः\par}

\vspace{1cm}

{\large Pages 58--114 of the B\={\i}jaga\d{n}ita\par}

\vspace{0.5cm}

{\normalsize\itshape To be continued in Part III (pages 115--170)\par}

\end{center}



\vfill



\begin{center}

\rule{0.5\textwidth}{0.4pt}\\[0.5em]

{\small\textsc{Bilingual Sanskrit-English Edition}\\

\textsc{Sanskrit text: Devanagari} $\parallel$ \textsc{English: IAST transliteration}\\[0.3em]

\textsc{Compiled from the Sanskrit edition of}\\

\textsc{Shri Bhaskaracharya}\\

\textsc{Circa 1150 CE}\par}

\end{center}



\end{document}

%% ========================================================================

